\chapter{Diagnóstico y análisis del edificio}

Este capítulo se desarrollará a partir de la información técnica del proyecto arquitectónico, los planos disponibles, las especificaciones constructivas y el modelo digital del edificio. Su finalidad es establecer el escenario base de desempeño ambiental y energético antes de introducir estrategias de optimización.

\section{Descripción del proyecto}

El proyecto corresponde al edificio de la Defensoría del Pueblo en la isla San Cristóbal, Galápagos. La descripción arquitectónica deberá incluir implantación, orientación, programa funcional, sistema constructivo, materialidad, configuración de envolvente, aperturas, protecciones solares y criterios de ventilación natural existentes.

\section{Condiciones climáticas del contexto}

El análisis climático considerará temperatura, humedad relativa, radiación solar, vientos dominantes y condiciones propias del contexto insular. Estos datos permitirán identificar oportunidades de control solar, ventilación natural, reducción de cargas térmicas y mejora del confort adaptativo.

\section{Escenario base de desempeño}

El escenario base corresponde al diseño actual del proyecto. Se evaluará su comportamiento frente a criterios de eficiencia energética y confort térmico, considerando demanda energética, ganancias térmicas, ventilación, iluminación y condiciones interiores proyectadas.

\section{Matriz de indicadores}

La matriz de indicadores integrará criterios provenientes de estándares como Minergie, WELL, CEELA, EDGE y LEED, junto con normativa local aplicable a Galápagos. Esta matriz permitirá evaluar el nivel de respuesta del proyecto frente a parámetros mínimos de sostenibilidad, eficiencia y confort.

\section{Oportunidades de mejora}

A partir del diagnóstico se identificarán oportunidades de mejora asociadas a orientación, protección solar, ventilación natural, aislamiento térmico, selección de materiales, reducción de huella de carbono y optimización del desempeño ambiental del edificio.
